\documentclass{article}
\usepackage[T1]{fontenc}
\usepackage{lmodern}
\usepackage{graphicx}
\usepackage{amsmath,amssymb}
\usepackage[a4paper,margin=2cm]{geometry}
\usepackage[absolute,overlay]{textpos}
\setlength{\TPHorizModule}{1mm}
\setlength{\TPVertModule}{1mm}
\usepackage[indonesian]{babel}
\usepackage{hyperref}
\setlength{\emergencystretch}{2em}
% unit: Halaman 2

\begin{document}

% === Halaman 2 (llm) ===

\section*{Materi Matematika untuk Kriptorafi:}

\begin{enumerate}
    \item Teori Bilangan
    \begin{itemize}
        \item Integer dan aritmetika modulo
        \item Algoritma euclidean
        \item Kekongruenan
        \item Relatif prima
        \item Balikan (invers) modulo
        \item Bilangan prima
    \end{itemize}
    \item Probabilitas dan Statistik
    \item Kompleksitas algoritma
    \item Teori Informasi
    \item Aljabar abstrak
\end{enumerate}

\vspace{10mm}

\begin{tabular}{|p{9cm}|}
\hline
No. 1 s/d 3 sudah dipelajari di dalam kuliah Matematika Diskrit dan Probabilitas dan Statistik \\
\hline
No 5 akan dibahas pada materi $ECC$ (Elliptic Curve Cryptography) \\
\hline
\end{tabular}

\end{document}
